\documentclass[11pt,twocolumn]{article} 
\usepackage{fullpage}


\begin{document}

\title{A writeup for my favorite class}
\author{Me Myself}
\maketitle 


\abstract{This writeup summarizes the design and findings of my first
  project in csci350, Spring09. }

\section{Overview}

This project acts as a miniature GIS and lets the user execute the
followig commands: 
\begin{verbatim}
gis>help
2drender grid.ass
3drender grid.asc
flowdir elev.asc dir.asc
flowaccu elev.asc dir.asc accu.asc
bye
gis>
\end{verbatim}

To start a new paragraph you need to start a new line. 

The \texttt{flowdir}  program computes the flow direction grid
corresponding to an elevation file. 


The \texttt{flowaccu} program computes the flow accumulation grid
corresponding to an elevation file and a flow direction file.


\section{Algorithm and analysis}

\paragraph{Flow direction. }
The flow direction computation works by
iterating in order over the grid, and computing, for each point
$(i,j)$, the neighbor of minimum height. A neighbor is considered a
point $(i+/-1, j+/-1)$.

{\bf Space:} My algorithm is very space-aware. Or is it
space-oblivious. 

{\bf Time: } And the running time is also outstanding.


\paragraph{Flow accumulation.} And here is my flow accumulation
approach. The flow accumlaion algorithm works as follows...

{\bf Space:} The array is $a_1, a_2, ..., a_n$. 

{\bf Time: } The running time is $O(n \log n)$ if $n > 3$ otherwise it
is $O(n^3)$. 




\section{Experimental evaluation}

To determine the performance in practice of my algorithms I performed
a careful experimental evaluation.

{\bf Platform.} The code was written in C, compiled with \texttt{gcc
  4.0.2, -O3}, 64 bits enable, and executed on a Mac dual core 2.xx
GHZ, 4GB RAM you check all these things, right? running Leopard.

{\bf Datasets: } The test datasets are real-life grid terrains ranging
from 1 million to 1,000 million points.  They are given in
Table~\ref{tbl:datasets} below.


\begin{table}[htp]
    \centering{
    \begin{tabular}{|l|r|}
      \hline 
      Dataset & nb. of points \\
      \hline 
      \hline 
      {Kaweah}        &      $1.6 \cdot 10^6$ \\
      {Sierra}        &      $9.5 \cdot 10^6$ \\
      {East-Coast USA (usadem2)} & $ 246 \cdot 10^6$ \\
      {Washington State}    & $1,066 \cdot 10^6$ \\
      \hline
    \end{tabular}
  }
  \caption{Size of terrain datasets.}
  \label{tbl:datasets}
\end{table}
      

The running times for \texttt{flowaccu} and \texttt{flowdir} are
summarized in Tables~\ref{tbl:rundir} and \ref{tbl:runaccu}. Don't
worry about the placement of tables cause Latex places the figures and
table the best it can. That's not always where you want them.

\begin{table}[htp]
    \centering{
    \begin{tabular}{|l|r|r|}
      \hline 
      Dataset & time & CPU\\
      \hline 
      \hline 
      {Kaweah}        &     ? & ? \\
      {Sierra}        &      ? & ? \\
      {East-Coast USA (usadem2)} & ? & ? \\
      {Washington State}    & ? & ? \\
      \hline
    \end{tabular}
  }
  \caption{Running times (seconds) and CPU-utilization for
    \texttt{flowdir}.}
  \label{tbl:rundir}
\end{table}


\begin{table}[htp]
    \centering{
    \begin{tabular}{|l|r|r|}
      \hline 
      Dataset & time & CPU\\
      \hline 
      \hline 
      {Kaweah}        &     ? & ? \\
      {Sierra}        &      ? & ? \\
      {East-Coast USA (usadem2)} & ? & ? \\
      {Washington State}    & ? & ? \\
      \hline
    \end{tabular}
  }
  \caption{Running times (seconds) and CPU-utilization for
    \texttt{flowaccu}.}
  \label{tbl:runaccu}
\end{table}

%%in case you want to put a figure
% \begin{figure}[hptb]
%   \centering
%   \includegraphics [width=0.45 \linewidth] {figures/kaweah2.pdf} 
%   \includegraphics [width=0.45 \linewidth]{figures/sierra.pdf}


\section{Conclusion}
The conclusion goes here.  

\end{document}


